\documentclass[11pt,a4paper]{article}
\usepackage[margin=25mm]{geometry}
\usepackage{amsmath,amssymb,amsthm,hyperref}
\setlength{\emergencystretch}{2em}
\newtheorem{lemma}{Lemma}
\newtheorem{theorem}[lemma]{Theorem}
\newcommand{\gen}[1]{\langle#1\rangle}
\title{A transitive four-point residual fixed set}
\author{Human Compression project: L044 continuation}
\date{}
\begin{document}\maketitle

Let $G=\gen{\rho_0,\ldots,\rho_{r-1}}\leq A_n$ be a transitive
string C-group whose noncommuting diagram is connected and whose every
maximal omission is intransitive. Suppose $G_0$ has orbits $\{a\}$ and
$B=\Omega\setminus\{a\}$, and write
\[
 \rho_0=(a\ b)\delta,\qquad X=\Omega\setminus\{a,b\},\qquad
 K=\gen{\rho_2,\ldots,\rho_{r-1}}.
\]
Thus $K$ fixes $a,b$ and centralizes the odd involution $\delta$ on $X$.
Let $\mathcal T$ be its set of $t$ two-point orbits, let $P$ be their
union, and let $F=\operatorname{Fix}_X(\delta)$. Assume
\[
 n=2r\geq14,\qquad |F|=4.
\]
Consequently $n=2t+6$, $t\geq5$ is odd, and $K$ has $t+1$ distinguished
generators. Put $\sigma=\rho_1$ and $R=\gen{\rho_3,\ldots,\rho_{r-1}}$.
The permutation $\sigma$ fixes $a$ and centralizes $R$, and
$G_0=\gen{\sigma,K}$ is transitive on $B$.

We use the accepted actual block-orbit deficit proof: deleting any
distinguished $K$ image on the $\delta$-blocks strictly refines the orbit
partition. In particular all those images are independent. We also retain
Whiston's independent-length bound $m-1$ in $S_m$ and its equality clause,
which identifies the generated group as $S_m$. The accepted maximal
involutory string theorem says that $m-1$ independent string involutions
on an odd $m\geq5$ form the ordinary Coxeter path; its degree-five finite
certificate remains a paid dependency. No quotient intersection property
or even-degree maximal-string classification is used below.

\begin{lemma}[Transport of a natural orbit]\label{transport}
Suppose $U\subseteq F$ is an $R$-orbit on which the image is the natural
$S_s$, $s\geq3$. If $\sigma(U)=O\subseteq P$, then $O$ and $O^\delta$
are disjoint $R$-orbits. Their union consists of exactly $s$ pairs in
$\mathcal T$. Every member of $R$ fixing $U$ pointwise fixes both orbits
pointwise, and therefore fixes those $s$ pair-blocks.
\end{lemma}
\begin{proof}
Both $\sigma$ and $\delta$ commute with $R$, so they transport its
orbits equivariantly. If $O^\delta=O$, the restriction of $\delta$ to $O$
centralizes the natural $S_s$ action. This centralizer is trivial: an
element fixing one point fixes every point, and the point stabilizer
$S_{s-1}$ fixes only that point when $s\geq3$. Thus $\delta$ would fix
$O$ pointwise, contrary to $O\subseteq P$. The two orbits are disjoint,
and $\delta$ pairs them. Equivariance also gives the final assertion.
\end{proof}

\begin{theorem}\label{main}
Under the stated hypotheses, $K$ is not transitive on $F$.
\end{theorem}
\begin{proof}
Suppose it is transitive. We distinguish the possible restrictions of
the distinguished involutions to the four points of $F$.

\medskip\noindent\textbf{A double transposition occurs on $F$.}
There are two distinguished labels whose restrictions generate a
transitive group on $F$. Indeed, two distinct double transpositions
already do so. If only one double transposition occurs, transitivity
requires a transposition joining its two pairs, and these two suffice.
Call these labels the extras and let $J$ be the other $t-1$ labels.

For each label in $J$, the full omission still contains the two extras
and is transitive on $F$. Therefore its pair-block orbit partition must
change. Its pair image cannot belong to the subgroup on the other pair
images. In particular the $t-1$ pair images in $J$ are independent.
Whiston makes their group $S_t$, and the accepted odd-degree extremal
string theorem makes them an ordinary path. Consecutive path marks
do not commute, so their indices form an interval in the full string.
Write its pair vertices in path order as $C_1,\ldots,C_t$.

Each extra has trivial pair image. To see this, omit any path mark.
The remaining path has precisely the two orbits separated by that cut.
If either extra joined these two orbits, the full omission would again
be transitive on the pair-blocks and on $F$, contradicting the actual
block-orbit deficit. Thus each extra preserves every path cut.
Their common setwise stabilizer fixes every $C_i$.

If $\rho_2$ is an extra, $R$ contains all of $J$. Its pair image is
transitive, so its point orbits in $P$ have size $t$ or $2t$: each such
orbit meets every two-point block in the same positive number of points.
All its other orbits have size at most four. Hence $P$ is precisely
the union of its orbits of size at least five, and $\sigma$ preserves $P$.
So does $K$, contradicting transitivity of $G_0$ on $B$.

Otherwise $J$ starts at $\rho_2$, and the two extras $\eta,\theta$ occur
consecutively at its far end, in that order. On $F$ they cannot both be
transpositions, as two transpositions do not act transitively on four
points. They cannot both be double transpositions either: these commute
on $F$, and their trivial pair images make their restrictions to $P$
products of pair flips, which commute too. They would commute globally,
disconnecting the full string. Exactly one extra is therefore a
transposition on $F$.

The outer extra $\theta$ commutes with every path mark. Its flip vector
is constant on the $t$ pairs. If it is odd on $F$, its evenness globally
forces it to flip every pair, since $t$ is odd. If it is even on $F$,
the inner extra $\eta$ is odd there. This inner extra commutes with all
but the last path mark, so its flip vector is constant on
$C_1,\ldots,C_{t-1}$. Since $t-1$ is even, its global evenness forces
a flip on $C_t$.

The group $R$ now has pair-block orbits $\{C_1\}$ and
$\{C_2,\ldots,C_t\}$. It is transitive on the points over the latter
orbit: in the first case $\theta$ flips both points in each block;
in the second $\eta$ flips $C_t$, and pair transitivity spreads this
within-block connection to every block. The resulting point orbit has
size $2(t-1)\geq8$. The extras make $F$ an $R$-orbit of size four.
All other orbits have size at most two, so $F$ is its unique four-point
orbit. Centralization forces $\sigma$ to preserve $F$, again impossible.
This argument includes $t=5$; a possible four-point row is eliminated
by the explicitly forced flip, not by an even-degree rigidity claim.

\medskip\noindent\textbf{All nonidentity restrictions to $F$ are transpositions.}
Their edge graph is connected. No vertex has three incident distinct
edges: the three corresponding transpositions would pairwise fail to
commute, whereas three distinct positions in a string cannot be pairwise
adjacent. The graph is therefore a path or a four-cycle. A four-cycle
is impossible, since its four noncommuting adjacency requirements would
embed a cycle in the path of generator positions. Thus the graph is a
three-edge path.

There is exactly one label on each edge. Choose a label on the middle
edge. Labels on both outside edges must occupy its two neighboring
positions. A second middle-edge label would have to neighbor both
those positions, and a second outside-edge label would have to occupy
one of the same already occupied neighboring positions. Neither is
possible. We therefore have exactly three consecutive labels
$\eta,\zeta,\theta$ acting as the ordinary $S_4$ path on $F$; all the
remaining $t-2$ labels, denoted $M$, fix $F$ pointwise.

The pair images of $M$ are independent. If one belonged to the subgroup
on the other pair images, its full block image would belong to that
subgroup too, because every $M$ label fixes $F$. This contradicts the
accepted block-image independence. Every subfamily remains independent.

First suppose $\rho_2$ belongs to $M$. Then $R$ is transitive on $F$
with natural image $S_4$. If $\sigma$ preserved $F$, so would $G_0$.
It must therefore send $F$ to a different $R$-orbit in $P$; neither
$a$ nor $b$ can belong to an orbit of size four. Lemma~\ref{transport}
gives four pair-blocks fixed by all the $t-3$ labels in
$M\setminus\{\rho_2\}$. Their independent pair images consequently
act on just $t-4$ points. Whiston bounds their number by $t-5$, a
contradiction. This also covers $t=5$, when that action is trivial.

It follows that the three $F$ labels are $\rho_2,\rho_3,\rho_4$.
The group $R$ has a natural three-point orbit $U\subset F$ and fixes
the fourth point $z$. If $\sigma(U)$ were different from $U$, it would
lie in $P$. Lemma~\ref{transport} would give three pair-blocks fixed
by all $t-2$ labels of $M$, contradicting the bound $t-4$ on the
remaining $t-3$ pair-blocks. Thus $\sigma$ preserves $U$.

The point $\sigma(z)$ cannot be in $U$, by bijectivity, or equal $a$,
since $\sigma$ fixes $a$. If it were $z$, then $G_0$ would preserve
$F$; if it were $b$, the involution $\sigma$ would interchange $z,b$
and $G_0$ would preserve $F\cup\{b\}$. Therefore
$x=\sigma(z)$ lies in $P$. Since $R$ fixes $z$ and commutes with
$\sigma$ and $\delta$, it fixes both $x$ and $x^\delta$.
Write their pair as $C$.

All $t-2$ independent pair images of $M$ fix $C$. Equality in Whiston's
bound now makes their image the full $S_{t-1}$ on the other pairs.
The generator $\zeta=\rho_3$ commutes with every member of $M$ and
fixes $C$ pointwise, since it belongs to $R$. Its pair image centralizes
this $S_{t-1}$ and fixes $C$, so is trivial. Its flip vector is constant
on the other $t-1$ pairs. That number is even; hence $\zeta$ is even
on $P$, while its restriction to $F$ is a transposition. It fixes $a,b$,
so is odd globally, the final contradiction.
\end{proof}

\paragraph{Attribution and exact boundary.}
The singleton-end strategy comes from Cameron--Fernandes--Leemans--Mixer,
\emph{Highest rank of a polytope for $A_n$}, arXiv:1605.09173v1,
Section~5, especially Proposition~5.29 on physical page~38 and the
subsequent intransitive-tail branches. The accepted residual deficit is
due to the preceding project continuation; Whiston and the full
odd-degree extremal-string proof, including its old degree-five
certificate, remain retained inputs. The two-generator fixed-set split,
orbit transport, and final parity argument are the present addition.

This note excludes only the stated transitive four-point residual case.
It does not assume $G_0=A_B$, any smaller-degree alternating induction,
or faithfulness/IP of a pair quotient. It does not exclude additional
head configurations by convention. Together with the separately reviewed
small-orbit exclusions it can close $|F|=4$ in that precise endpoint
branch, but larger fixed sets, other source branches, rigidity and table
obligations remain separate. No full L044 recovery, original count,
compression factor, conjecture answer, priority or new formalization is
claimed. There is no new finite premise or numerical search.
\end{document}
